\hypearticle{Pesquisas eleitorais: Como ouvir poucos e retratar muitos?}
{Gustavo Neves Moraes}
{

\hspace*{1.5em}
Em períodos eleitorais, é comum ver pesquisas que entrevistam cerca de 2 mil pessoas
para descrever a intenção de voto de mais de 150 milhões de eleitores.
A desconfiança acerca dessa desproporção é natural,
mas ela se apoia em uma intuição equivocada de como é feita a seleção.
A precisão de uma pesquisa depende muito mais de como a amostra é escolhida
do que da fração do eleitorado que foi ouvida.

\subtitle{Por que uma amostra pequena pode bastar}

\hspace*{1.5em}
Para entender de forma simplificada, vamos usar uma alusão bastante famosa: 
para saber se uma sopa está no ponto, basta provar uma colherada, desde que a panela esteja bem misturada.
O tamanho da panela não muda isso,
pois o que importa é que a colherada represente o conjunto.

\hspace*{1.5em}
Em uma pesquisa eleitoral, a mistura da sopa corresponde à seleção aleatória, isto é,
quando todo eleitor tem uma chance conhecida de ser escolhido,
a amostra tende a reproduzir as proporções do eleitorado.
Ainda assim, ela não será idêntica à população,
já que cada sorteio produz um grupo ligeiramente diferente do outro.
Essa diferença é meramente fruto do acaso e, por seguir regras matemáticas conhecidas,
pode ser medida e informada junto com o resultado, na forma da margem de erro.

\subtitle{Quanto a amostra pode errar}

\hspace*{1.5em}
Seja $p$ a proporção de eleitores que apoiam determinado candidato.
A proporção observada na amostra, $\hat{p}$, varia de uma amostra para outra,
e essa variação é medida pelo erro padrão.
Para um nível de confiança de 95\%, a margem de erro é aproximadamente:

\[
\text{margem} \approx 1{,}96 \sqrt{\frac{p(1-p)}{n}}
\]

\hspace*{1.5em}
O valor 1,96 vem da distribuição normal padrão.
Para amostras grandes, $\hat{p}$ se distribui aproximadamente como uma normal em torno de $p$,
e, nessa distribuição, cerca de 95\% da probabilidade fica a até 1,96 erros padrão da média.
A margem de erro, portanto, é esse número de erros padrão
multiplicado pelo erro padrão de $\hat{p}$, que é a raiz da fórmula.
Para outros níveis de confiança, o valor muda de acordo com a Tabela da Distribuição Normal (por
exemplo, valores famosos são 1,645 para 90\% e 2,576 para 99\%).

\hspace*{1.5em}
O tamanho da população não aparece na fórmula,
pois, para populações grandes, a precisão é determinada por $n$.
O produto $p(1-p)$ é máximo quando $p = 0{,}5$, e é esse caso que costuma servir de referência.
Como $n$ está dentro da raiz, reduzir a margem pela metade exige quadruplicar a amostra,
o que explica por que a maioria das pesquisas se limita a alguns milhares de entrevistados.

\subtitle{Um exemplo numérico}

\hspace*{1.5em}
Considere uma pesquisa com $n = 2000$ entrevistados em que um candidato aparece com 48\%.
A margem de erro é $1{,}96\sqrt{0{,}48 \cdot 0{,}52 / 2000} \approx 2{,}2$ pontos percentuais,
o que leva ao intervalo de 45,8\% a 50,2\%.
Isso não quer dizer que o valor verdadeiro esteja nesse intervalo com 95\% de probabilidade.
O que se afirma é que, se a pesquisa fosse repetida muitas vezes com amostras distintas,
cerca de 95\% dos intervalos construídos dessa forma conteriam o valor verdadeiro.

\subtitle{O que a margem de erro não cobre}

\hspace*{1.5em}
Vale ressaltar que a margem de erro mede apenas a variação causada pelo sorteio da amostra.
Ou seja, ela não captura problemas de outra natureza,
como a recusa de parte dos sorteados em responder, quando os que aceitam diferem dos que recusam,
ou o uso de amostragem por cotas, adotada por muitos institutos,
em que os entrevistados são escolhidos para reproduzir perfis de sexo, idade, escolaridade e renda,
e não sorteados individualmente.
Nesses casos, a margem divulgada é uma aproximação, e o erro real pode ser maior.
Além disso, a pesquisa descreve a intenção de voto no momento da coleta e não prevê o resultado,
pois eleitores indecisos podem mudar de opinião até a votação.

\originalsource{https://www.pewresearch.org/short-reads/2016/09/08/understanding-the-margin-of-error-in-election-polls/}

\vspace{0.1em}
{\color{HypePurple2} Quer se aprofundar mais?}

\href{https://www.youtube.com/shorts/cEq4mvdK10o}{{\color{HypeOrange}Confira este vídeo curto}} que instigou a escrita do tema. \\
\href{https://www.youtube.com/watch?v=85NM1-FQrn0}{{\color{HypeOrange}Confira este vídeo do Khan Academy}} sobre como pequenas amostras podem representar a opinião de milhões de eleitores.

}